\section{两大阵营与三大模型：交替领先为什么会成为常态}

前两章讨论投入和收入，本章进入竞争格局。这里的问题是：如果大家都继续投钱，产业会围绕哪些主线组织起来？广密把 AI War 的硬件生态简化成两条主线：Nvidia GPU 生态和 Google TPU 生态；把应用/模型层主线简化成 OpenAI、Anthropic、Google Gemini 的三强交替领先。这个简化很有用，因为它避免把每个季度的新榜单都当成终局，也提醒我们把“谁短期领先”放回更长期的生态竞争中理解。

\subsection{Nvidia GPU vs Google TPU}

本节先从底层硬件生态讲起，因为模型公司的能力、成本和迭代速度都要落在芯片、软件栈和云供给上。理解 GPU 与 TPU 的差异，才能理解为什么同样是模型竞争，Google 和 Nvidia 代表的是两种完全不同的产业组织方式。

TPU 是 Google 自研的 Tensor Processing Unit，面向机器学习训练和推理优化；GPU 是 Nvidia 主导的通用并行计算生态，围绕 CUDA、开发者、云服务、服务器、网络和供应链形成了强大网络效应。Google 的优势是端到端：模型、芯片、云、搜索、Android、YouTube 和 Workspace 都在一个公司内；Nvidia 的优势是开放生态：OpenAI、Anthropic、xAI、Meta、云厂商和创业公司都可以围绕 GPU 形成外部创新网络。

广密把 Google 比作 AI 时代的苹果，把 Nvidia 比作 AI 时代的安卓。这不是说二者完全等价，而是指出两种组织方式：垂直整合与开放生态。前者在产品和成本上可能形成闭环，后者在开发者、多公司创新和供给扩张上更有弹性。

\begin{knowledgebox}{术语消化：GPU、TPU 与 CUDA}
\begin{tabularx}{\textwidth}{p{0.18\textwidth}p{0.32\textwidth}X}
\toprule
术语 & 解决的问题 & 本期关系 \\
\midrule
GPU & 用大规模并行计算加速矩阵运算和深度学习 & Nvidia 生态的核心，支撑 OpenAI/Anthropic 等外部模型公司。 \\
TPU & Google 为机器学习定制的加速芯片 & 支撑 Google 从芯片到 Gemini 的垂直闭环。 \\
CUDA & Nvidia 的并行计算软件生态 & 让 GPU 不只是硬件，而是开发者和工程工具链。 \\
生态位 & 公司在产业链里的稳定位置 & Google 做全栈，Nvidia 做底层生态，模型公司做能力和产品。 \\
\bottomrule
\end{tabularx}
\end{knowledgebox}

\subsection{GPT、Claude、Gemini 的交替领先}

三大模型的交替领先，是本期最重要的格局判断之一。广密认为，在现有 Pre-training 加 RL/Post-training 范式下，顶尖团队的方法、人才和信息流动很快，单家公司很难长期拉开巨大差距。于是竞争从“谁一次性碾压”变成“谁在某个战略 Beta 上端到端优化得更彻底”。

\begin{figure}[H]
\centering
\includegraphics[width=0.96\textwidth]{top-model-rotation.png}
\caption{三大模型交替领先：GPT、Claude、Gemini 的战略重点不同。自制概念图，依据 00:17:48--00:25:40 对谈内容整理。}
\end{figure}

\begin{knowledgebox}{读图：交替领先不是没有分化}
GPT 更强在 C 端产品、品牌和个人助理心智；Claude 更强在 Coding、企业和知识工作者体验；Gemini 更强在多模态、Google 生态和 Android/搜索分发。交替领先说明通用能力难以长期拉开，但分化说明每家公司会围绕自己的数据和产品目标优化。
\end{knowledgebox}

\begin{warningbox}{Benchmark SOTA 的误读}
SOTA 原意是 state of the art，即当前最先进水平。但本期强调：榜单 SOTA 不等于商业 SOTA。真正重要的指标可能是营收、客户、日活、留存、tokens 输出、任务成功率和用户愿意交付多少真实工作。
\end{warningbox}

\subsection{基础模型公司为什么像综合电商}

广密提出一个“偷懒但有用”的类比：基础模型公司像十年前的综合电商，Scale SKU 类比 Scale Data。综合电商从书扩展到全品类时，可以复用仓储、物流、支付、流量和用户关系；基础模型公司从聊天扩展到 Coding、Office、投研、医疗、设计和多模态时，也可以复用模型架构、数据管线、评测系统、RL 流程和用户入口。

\begin{figure}[H]
\centering
\includegraphics[width=0.94\textwidth]{foundation-model-as-ecommerce.png}
\caption{基础模型 = 综合电商：Scale SKU 类比 Scale Data。自制概念图，依据 00:25:40--00:27:40 对谈内容整理。}
\end{figure}

\begin{knowledgebox}{读图：类比的重点是“复用基础设施”}
综合电商能从图书扩到全品类，不只是因为想卖更多东西，而是仓配、支付、推荐和用户信任可以复用。基础模型公司能从聊天扩到 Coding、Office 和行业 Agent，也不是简单加功能，而是数据采集、训练、评测和产品反馈管线可以复用。
\end{knowledgebox}

这个类比对创业者有两个提醒。第一，垂直场景不一定没有价值，但要知道自己是在做“垂直电商”还是在做“小红书”。垂直电商可能赚钱，但平台上限受限；小红书在综合电商之外构建内容资产和社区心智，才有不同价值。第二，模型公司向下做数据的效率更高，创业公司必须找到基础模型公司不容易获得的差异化数据、工作流或用户关系。

\subsection{本章小结}

本章把竞争格局拆成硬件生态、模型三强和平台类比。交替领先说明现有范式下技术差距难以永久拉开；战略 Beta 和数据闭环则决定每家公司在哪里建立特色。基础模型像综合电商，意味着“横向扩品类”的能力很强，创业公司必须避开被轻易 Scale Data 的位置。

